% ==============================================================================
% CẤU HÌNH BỔ SUNG: tài liệu tham khảo, liên kết, macro của đề tài
% ==============================================================================

% Tài liệu tham khảo dạng [n] (IEEE) qua BibLaTeX/Biber
\usepackage[
    backend=biber,
    style=ieee,
    sorting=none,
    doi=true,
    isbn=false,
    url=true
]{biblatex}
\addbibresource{references.bib}
\setlength{\bibitemsep}{3pt}

% Liên kết nội bộ (không tô màu, giống bản in Word)
\usepackage{hyperref}
\hypersetup{
    hidelinks,
    pdfauthor={Nhóm 7 - Trần Trọng Hoàn, Nông Hữu Thịnh, Nguyễn Quang Minh, Vũ Ngọc Quý, Trần Văn Dưỡng},
    pdftitle={\projectTitle},
    pdfsubject={Báo cáo dự án học phần \courseName},
    pdfkeywords={Nhận diện biển báo giao thông, thị giác máy tính, học sâu}
}

% Macro tiện ích
\newcommand{\code}[1]{\texttt{#1}}
\newcommand{\abbr}[2]{\textbf{#1} (#2)}

% ------------------------------------------------------------------------------
% Bổ sung riêng cho đề tài
% ------------------------------------------------------------------------------
\usetikzlibrary{positioning,arrows.meta,fit,calc,shapes.geometric,backgrounds}
\usepackage{bm}
\usepackage{float}
\usepackage{pifont}
\usepackage{xurl}
\setcounter{biburlnumpenalty}{100}
\setcounter{biburllcpenalty}{100}
\setcounter{biburlucpenalty}{100}
% Ghi chú ngắn dưới bảng / hình (như báo cáo môn Học sâu)
\newcommand{\tablenote}[1]{\par\vspace{3pt}{\centering\small #1\par}}
\newcommand{\fignote}[1]{\par\vspace{2pt}{\centering\small #1\par}}
